\section{SDE/ODE 视角下的引导生成}

\subsection{SDE 表述中的引导}

扩散模型的反向过程可以等价地表述为连续时间的随机微分方程（SDE）。在 SDE 表述下，反向过程的时间域离散化形式为：

$$x_{t-\Delta t} = x_t + f(x_t, t)\Delta t + g(t)\Delta W$$

其中 $f$ 包含分数函数，$g$ 是噪声系数。将 DPS 引导引入 SDE 表述同样是直接的：在每一步计算 $x_{t-\Delta t}$ 后，根据条件梯度进行更新。

\begin{importantbox}{SDE 与 DDPM 的等价性}
DDPM 的离散反向步骤可以视为对连续 SDE 的时间域离散化。两种表述完全等价——SDE 表述提供了更自然的连续时间框架，而 DDPM 是其在离散时间步上的特例。DPS 引导可以在两种表述中一致地应用。
\end{importantbox}

\subsection{ODE 表述的可行性}

除了 SDE 表述外，扩散模型的反向过程也可以用概率流 ODE 来描述。概率流 ODE 的形式为：

$$\frac{dx}{dt} = f(x, t) - \frac{1}{2}g(t)^2 \nabla_x \log p_t(x)$$

其特点是没有随机噪声项，给定初始条件 $x_T$ 的轨迹是确定性的。课堂讨论表明：DPS 的引导思路完全可以应用到 ODE 表述中，只需在每步求解后添加引导梯度更新。

\begin{knowledgebox}{三种表述的对比}
\begin{center}
\begin{tabular}{lccc}
\toprule
特性 & DDPM & SDE & ODE \\
\midrule
时间域 & 离散 & 连续 & 连续 \\
随机性 & 有 & 有 & 无 \\
采样质量 & 高 & 高 & 高 \\
采样速度 & 慢（需多步） & 慢 & 可加速 \\
DPS 引导 & 支持 & 支持 & 支持 \\
理论基础 & 变分推断 & 随机分析 & 概率流 \\
\bottomrule
\end{tabular}
\end{center}
三种表述在引导方面面临相同的实际问题：Tweedie 近似的精度和引导强度的权衡。
\end{knowledgebox}

\subsection{时间步数与引导强度的权衡}

\begin{warningbox}{更多时间步 vs. 更大引导权重}
推理时引导存在一个核心权衡：
\begin{itemize}
    \item \textbf{增加时间步数}：更多步意味着更多次注入引导梯度，引导效果更好，但生成速度变慢。
    \item \textbf{增大引导权重}：减少时间步但增大每步的引导强度，会导致中间样本偏离训练分布，使得噪声预测网络的预测变得不准确，最终输出变得不真实。
\end{itemize}
这是推理时引导方法的核心难题：输出要么不满足给定条件，要么满足条件但变得不真实。如何在真实性与条件满足之间取得平衡，是调参的关键。
\end{warningbox}

如果希望同时获得加速（少步数）和良好引导效果，可能需要在训练阶段就引入条件信息。但这就失去了推理时引导``无需额外训练''的优势。

讲者指出，这一权衡在实践中表现为：当引导权重过大时，中间样本 $x_t$ 会被推到训练数据分布之外的区域。在这些分布外区域，噪声预测网络 $\epsilon_\theta$ 从未见过类似的输入，因此其预测会变得不准确。这种不准确的预测会在后续步骤中累积，导致最终输出虽然可能满足给定条件，但在视觉上变得不真实、产生伪影。

\begin{knowledgebox}{实用调参建议}
根据课堂讨论，以下经验法则有助于在实践中取得较好的结果：
\begin{itemize}
    \item 使用较多的去噪步数（如 1000 步），让每步引导的扰动较小
    \item 引导权重可以随时间步变化：在早期（高噪声）步使用较小的权重，在后期（低噪声）步逐渐增大
    \item 同时尝试 SDE 和 ODE 两种采样器，比较结果
    \item 对于同一个条件，多次采样取最优（即 best-of-N 策略）
\end{itemize}
\end{knowledgebox}

\subsection{本章小结}

DPS 引导可以在 DDPM、SDE 和 ODE 三种表述中一致应用。实践中需要仔细调整时间步数和引导权重，以在条件满足度和输出真实性之间取得平衡。

